\section{Onde aprofundar}

\subsection*{Livros}
\citetitle{livro:pemberton1982} \parencite{livro:pemberton1982}, com o
texto completo publicado pelo autor, é a melhor leitura sobre o p-code:
compilador e interpretador comentados. \citetitle{livro:goldberg1983}
\parencite{livro:goldberg1983}, capítulos~26 a~30, especifica a VM do
Smalltalk-80.

\subsection*{Papers}
\small
\begin{longtable}{@{}p{5.2cm} p{7cm}@{}}
\toprule
Fonte & Para quê \\
\midrule \endhead
\textcite{paper:aycock-history-jit} & Survey de quarenta anos de JIT,
  de 1960 a 2000: o mapa para o resto da lista \\
\textcite{paper:self-pics} & Inline caches polimórficos; curto e com PDF
  aberto \\
\textcite{paper:self-deoptimization} & De onde vem a desotimização \\
\textcite{paper:holzle-thesis-adaptive-optimization} & Otimização
  adaptativa completa: o desenho da HotSpot, em 1994 \\
\textcite{paper:hotspot-server-compiler} & O C2 descrito por quem o
  fez; PDF aberto na USENIX \\
\textcite{paper:clr-generics} & Generics com tipos exatos em runtime, o
  contraponto ao erasure \\
\textcite{paper:vm-showdown} & Pilha contra registradores, com dados \\
\bottomrule
\end{longtable}
\normalsize

\subsection*{Documentação}
\citetitle{doc:ecma335} (o padrão da CLI, com download do PDF na página
da Ecma), \citetitle{doc:android-dalvik-bytecode},
\citetitle{doc:android-art} e \citetitle{doc:android-art-jit}.

\subsection*{De quem fez o Java}
Nomes, bios e mais obras de cada pessoa no apêndice~\ref{ap:pessoas}.
\begin{itemize}
\obra*{paper:java-language-environment}{O white paper de 1996: o que o
  Java herdou e o que recusou das linguagens e VMs deste capítulo.}
\obra*{paper:feel-of-java}{Gosling sobre as decisões de design, em
  primeira pessoa.}
\obra*{livro:lindholm2014}{O item~1.1, \enquote{A Bit of History}, conta
  de onde veio a JVM.}
\obra*{video:java-story}{O documentário oficial, com quem construiu o
  Java.}
\end{itemize}

\subsection*{Pesquisa acadêmica}
Itens que completam os papers acima. Acesso e links no
apêndice~\ref{ap:pesquisa}.
\begin{itemize}
\pesquisa{paper:pascal-p-code-generation}{O compilador Pascal-P de Zurique
  e a máquina de pilha hipotética que a UCSD adaptou no p-System.}
\pesquisa{paper:self-power-of-simplicity}{A linguagem Self: protótipos,
  slots e mensagens, o começo da linha Self, Strongtalk, HotSpot.}
\pesquisa{paper:self-type-feedback}{Type feedback, a ideia central dos JITs
  modernos.}
\pesquisa{paper:jalapeno-adaptive-optimization}{A otimização adaptativa da
  JVM de pesquisa da IBM, que virou a Jikes RVM.}
\pesquisa{paper:lua50-implementation}{A primeira VM de registradores de uso
  amplo, feita na PUC-Rio.}
\pesquisa*{paper:hazan-maquina-registradores-java}{Pilha ou registradores
  para Java, com o olhar de dispositivos móveis (PUC-Rio).}
\pesquisa{paper:jit-value-specialization}{Pesquisa da UFMG em JIT
  especulativo: especializar funções JavaScript para os valores reais.}
\pesquisa{paper:isamap-dbt}{Tradução dinâmica de binários do IC-Unicamp,
  parente do JIT: traduzir e guardar código em execução.}
\end{itemize}
